% !TEX program = xelatex
% Самодостаточный конспект. Компилировать XeLaTeX/LuaLaTeX (tectonic 05-nesobstvennye-parametr-eyler.tex).
\documentclass[11pt,a4paper]{article}
\newcommand{\coursename}{Математический анализ}
\newcommand{\rightheader}{§5. Несобственные интегралы. Эйлеровы интегралы}
\newcommand{\doctitle}{Несобственные интегралы. Интегралы с параметром. Эйлеровы интегралы}
% ==== Общая преамбула конспектов (собирается в каждый .tex как самодостаточная) ====
% Перед этим файлом должны быть определены: \documentclass и макросы
%   \coursename  — название курса (левый колонтитул, подзаголовок),
%   \rightheader — правый колонтитул (напр. «§2. Рациональные дроби»),
%   \doctitle    — заголовок раздела на титуле.
\usepackage{fontspec}
\setmainfont{cmunrm.otf}[
  BoldFont       = cmunbx.otf,
  ItalicFont     = cmunti.otf,
  BoldItalicFont = cmunbi.otf]
\usepackage{polyglossia}
\setdefaultlanguage{russian}
\setotherlanguage{english}

\usepackage{amsmath,amssymb,amsthm,mathtools}
\usepackage[a4paper,margin=2.2cm]{geometry}
\usepackage{enumitem}
\usepackage{graphicx}
\usepackage{array}
\usepackage{xcolor}
\definecolor{accent}{HTML}{1F6FEB}
\definecolor{rulecol}{HTML}{D0D7DE}
\usepackage[colorlinks=true,linkcolor=accent,urlcolor=accent,citecolor=accent]{hyperref}
\usepackage{titlesec}
\usepackage{fancyhdr}

% ----- Колонтитулы -----
\pagestyle{fancy}
\fancyhf{}
\fancyhead[L]{\small\itshape \coursename}
\fancyhead[R]{\small\itshape \rightheader}
\fancyfoot[C]{\thepage}
\renewcommand{\headrulewidth}{0.4pt}
\renewcommand{\headrule}{\hbox to\headwidth{\color{rulecol}\leaders\hrule height \headrulewidth\hfill}}

% ----- Заголовки -----
\titleformat{\section}{\large\bfseries\color{accent}}{\thesection.}{0.6em}{}
\titleformat{\subsection}{\bfseries}{\thesubsection.}{0.5em}{}

% ----- Теоремные окружения (стиль конспекта) -----
\theoremstyle{definition}
\newtheorem{definition}{Определение}[section]
\newtheorem{example}[definition]{Пример}
\newtheorem{remark}[definition]{Замечание}
\theoremstyle{plain}
\newtheorem{theorem}[definition]{Теорема}
\newtheorem{lemma}[definition]{Лемма}
\newtheorem{corollary}[definition]{Следствие}
\newtheorem{property}[definition]{Свойство}
\newtheorem{criterion}[definition]{Критерий}
\renewcommand{\proofname}{Доказательство}

% ----- Операторы и сокращения -----
\DeclareMathOperator{\tg}{tg}
\DeclareMathOperator{\ctg}{ctg}
\DeclareMathOperator{\arctg}{arctg}
\DeclareMathOperator{\arcctg}{arcctg}
\DeclareMathOperator{\sh}{sh}
\DeclareMathOperator{\ch}{ch}
\DeclareMathOperator{\Th}{th}
\DeclareMathOperator{\cth}{cth}
\DeclareMathOperator{\Const}{const}
\DeclareMathOperator{\sign}{sign}
\DeclareMathOperator{\rang}{rang}
\DeclareMathOperator{\diam}{diam}
\DeclareMathOperator{\Span}{span}
\DeclareMathOperator{\Ker}{Ker}
\DeclareMathOperator{\Img}{Im}
\DeclareMathOperator{\tr}{tr}
\newcommand{\R}{\mathbb{R}}
\newcommand{\C}{\mathbb{C}}
\newcommand{\N}{\mathbb{N}}
\newcommand{\Z}{\mathbb{Z}}
\newcommand{\Q}{\mathbb{Q}}
\newcommand{\dx}{\,dx}
\newcommand{\dt}{\,dt}
\newcommand{\du}{\,du}
\newcommand{\eps}{\varepsilon}
\renewcommand{\Re}{\operatorname{Re}}
\renewcommand{\Im}{\operatorname{Im}}
\renewcommand{\le}{\leqslant}
\renewcommand{\ge}{\geqslant}
\newcommand{\scal}[2]{\left( #1,\,#2\right)}
\DeclarePairedDelimiter{\abs}{\lvert}{\rvert}
\DeclarePairedDelimiter{\norm}{\lVert}{\rVert}

\begin{document}
\begin{center}
  {\LARGE\bfseries \doctitle}\\[4pt]
  {\large\itshape \coursename}
\end{center}
\vspace{6pt}

\section{Несобственные интегралы первого рода}

Определённый интеграл Римана предполагает конечный промежуток и ограниченную функцию. Снятие
первого ограничения приводит к интегралам по бесконечному промежутку.

\begin{definition}[несобственный интеграл I рода]
Пусть $f$ непрерывна на $[a,+\infty)$ и интегрируема на каждом $[a,\xi]$. \emph{Несобственным
интегралом} от $f$ по $[a,+\infty)$ называют предел
\[
  \int_a^{+\infty}f(x)\dx=\lim_{\xi\to+\infty}\int_a^{\xi}f(x)\dx .
\]
Интеграл \emph{сходится}, если этот предел конечен; иначе (предел бесконечен или не существует) —
\emph{расходится}. Если интеграл сходится, $f$ называют интегрируемой в несобственном смысле.
\end{definition}

Аналогично определяют $\displaystyle\int_{-\infty}^{a}f=\lim_{\xi\to-\infty}\int_\xi^a f$, а
интеграл по всей оси полагают сходящимся, когда сходятся оба слагаемых:
\[
  \int_{-\infty}^{+\infty}f=\int_{-\infty}^{a}f+\int_a^{+\infty}f .
\]

\begin{example}[эталонный интеграл]
\[
  \int_1^{+\infty}\frac{dx}{x^\alpha}=\lim_{\xi\to+\infty}\frac{\xi^{1-\alpha}-1}{1-\alpha}
   \ (\alpha\ne1),\qquad
  \int_1^{+\infty}\frac{dx}{x}=\lim_{\xi\to+\infty}\ln\xi=+\infty .
\]
Отсюда $\displaystyle\int_1^{+\infty}\frac{dx}{x^\alpha}$ \textbf{сходится при $\alpha>1$ и
расходится при $\alpha\le1$}.
\end{example}

\section{Несобственные интегралы второго рода}

Снятие второго ограничения (ограниченности) даёт интегралы от неограниченных функций.

\begin{definition}[несобственный интеграл II рода]
Пусть $f$ непрерывна на $[a,b)$ и неограниченна в окрестности $b$ (особая точка). Полагают
\[
  \int_a^{b}f(x)\dx=\lim_{\eps\to0+}\int_a^{b-\eps}f(x)\dx .
\]
При особенности в точке $a$ — симметрично; при внутренней особой точке $c$ интеграл разбивают на
$\int_a^c+\int_c^b$ и требуют сходимости обоих.
\end{definition}

\begin{example}[эталонный интеграл]
\[
  \int_0^{1}\frac{dx}{x^\alpha}=\lim_{\eps\to0+}\frac{1-\eps^{1-\alpha}}{1-\alpha}\ (\alpha\ne1),
  \qquad \int_0^{1}\frac{dx}{x}=\lim_{\eps\to0+}(-\ln\eps)=+\infty .
\]
Значит, $\displaystyle\int_0^{1}\frac{dx}{x^\alpha}$ \textbf{сходится при $\alpha<1$ и расходится
при $\alpha\ge1$} — условие, противоположное случаю I рода.
\end{example}

\section{Свойства несобственных интегралов}

Везде далее $f$ определена на $[a,b)$, $b$ конечно или $b=+\infty$, и $f\in\mathcal R[a,\xi]$ при
любом $\xi<b$.

\begin{property}[линейность]
Если сходятся $\int_a^b f$ и $\int_a^b g$, то при любых $\lambda,\mu\in\R$ сходится и
$\displaystyle\int_a^b(\lambda f+\mu g)=\lambda\!\int_a^b f+\mu\!\int_a^b g$.
\end{property}

\begin{property}[Ньютон–Лейбниц]
Если $f$ непрерывна на $[a,b)$ и $F$ — её первообразная, то
$\displaystyle\int_a^b f=F(b-0)-F(a)$, где $F(b-0)=\lim_{\xi\to b}F(\xi)$ (когда предел конечен).
\end{property}

\begin{property}[интегрирование по частям]
Если $u,v$ имеют непрерывные производные на $[a,b)$ и существует конечный предел
$\lim_{\xi\to b}u(\xi)v(\xi)=u(b-0)v(b-0)$, то из сходимости одного из интегралов
$\int_a^b u\,dv$, $\int_a^b v\,du$ следует сходимость другого и
\[
  \int_a^b u\,dv=uv\Big|_a^{b-0}-\int_a^b v\,du .
\]
\end{property}

\begin{property}[замена переменной]
Если $f$ непрерывна на $[a,b)$, а $x=\varphi(t)$ строго возрастает, непрерывно дифференцируема и
$\varphi(\alpha)=a$, $\varphi(\beta-0)=b$, то
$\displaystyle\int_a^b f(x)\dx=\int_\alpha^\beta f(\varphi(t))\varphi'(t)\dt$ (при условии, что
хотя бы один из интегралов сходится).
\end{property}

\section{Признаки сходимости неотрицательных функций}

\begin{theorem}[критерий ограниченности]
Пусть $f\ge0$ на $[a,b)$. Несобственный интеграл $\int_a^b f$ сходится тогда и только тогда, когда
функция $\Phi(\xi)=\int_a^\xi f$ ограничена сверху:
$\exists C\ \forall\xi\in[a,b):\ \int_a^\xi f\le C$.
\end{theorem}

\begin{proof}
При $f\ge0$ функция $\Phi(\xi)$ не убывает. Монотонная функция имеет конечный предел при
$\xi\to b$ тогда и только тогда, когда она ограничена сверху (теорема о пределе монотонной
функции); этот предел и есть значение интеграла.
\end{proof}

\begin{theorem}[признак сравнения]
Пусть $f,g$ непрерывны на $[a,b)$ и $0\le f(x)\le g(x)$. Тогда:
\begin{itemize}[leftmargin=*]
  \item из сходимости $\int_a^b g$ следует сходимость $\int_a^b f$;
  \item из расходимости $\int_a^b f$ следует расходимость $\int_a^b g$.
\end{itemize}
\end{theorem}

\begin{proof}
Из $f\le g$ интегрированием: $\int_a^\xi f\le\int_a^\xi g$. Если $\int_a^b g$ сходится, то правая
часть ограничена сверху своим пределом, значит и $\int_a^\xi f$ ограничена — по критерию
ограниченности $\int_a^b f$ сходится. Вторая часть — контрапозиция первой.
\end{proof}

\begin{corollary}[предельный признак]
Если $f,g>0$ и $\displaystyle\lim_{x\to b}\frac{f(x)}{g(x)}=K$, $0<K<+\infty$, то интегралы
$\int_a^b f$ и $\int_a^b g$ сходятся или расходятся одновременно. Сравнение обычно ведут с
эталоном $1/x^\alpha$.
\end{corollary}

\begin{remark}[абсолютная сходимость]
Интеграл называют \emph{абсолютно сходящимся}, если сходится $\int_a^b\abs f$; абсолютная
сходимость влечёт сходимость. Сходимость, не являющаяся абсолютной, называется \emph{условной}
(пример — $\int_1^{+\infty}\frac{\sin x}{x}\dx$).
\end{remark}

\section{Интегралы, зависящие от параметра}

\begin{definition}
Пусть $f(x,y)$ определена на $[a,b]\times[c,d]$. Функцию
$\displaystyle I(y)=\int_a^b f(x,y)\dx$ называют интегралом, зависящим от параметра $y$.
\end{definition}

\begin{theorem}[непрерывность]
Если $f$ непрерывна на $[a,b]\times[c,d]$, то $I(y)$ непрерывна на $[c,d]$; в частности
\[
  \lim_{y\to y_0}\int_a^b f(x,y)\dx=\int_a^b\lim_{y\to y_0}f(x,y)\dx .
\]
\end{theorem}

\begin{theorem}[правило Лейбница]
Если $f$ и $f'_y$ непрерывны на $[a,b]\times[c,d]$, то $I(y)$ дифференцируема и
\[
  I'(y)=\frac{d}{dy}\int_a^b f(x,y)\dx=\int_a^b f'_y(x,y)\dx .
\]
\end{theorem}

\begin{theorem}[интегрирование по параметру, Фубини]
Если $f$ непрерывна на $[a,b]\times[c,d]$, то
\[
  \int_c^d\!\Bigl(\int_a^b f(x,y)\dx\Bigr)dy=\int_a^b\!\Bigl(\int_c^d f(x,y)\,dy\Bigr)\dx .
\]
\end{theorem}

\begin{example}[интеграл Дирихле]
Покажем $\displaystyle\int_0^{+\infty}\frac{\sin x}{x}\dx=\frac\pi2$. Введём
$\displaystyle I(y)=\int_0^{+\infty}e^{-xy}\frac{\sin x}{x}\dx$, $y\ge0$. Дифференцируя по
параметру (и дважды интегрируя по частям),
\[
  I'(y)=-\int_0^{+\infty}e^{-xy}\sin x\dx=-\frac{1}{1+y^2}.
\]
Значит $I(y)=-\arctg y+C$. Так как $I(y)\to0$ при $y\to+\infty$ (множитель $e^{-xy}\to0$), имеем
$0=-\frac\pi2+C$, $C=\frac\pi2$, и
\[
  \int_0^{+\infty}\frac{\sin x}{x}\dx=I(0)=\frac\pi2 .
\]
\end{example}

\begin{example}[дифференцирование по параметру]
$\displaystyle I(a)=\int_0^{\pi/2}\ln\bigl(a^2\cos^2x+b^2\sin^2x\bigr)\dx$, $a,b>0$. Тогда
\[
  I'(a)=\int_0^{\pi/2}\frac{2a\cos^2x}{a^2\cos^2x+b^2\sin^2x}\dx
   \overset{u=\tg x}{=\!=\!=}2\int_0^{\infty}\frac{a\,du}{(a^2+b^2u^2)(1+u^2)}=\frac{\pi}{a+b}.
\]
Интегрируя, $I(a)=\pi\ln(a+b)+C$. В случае $a=b$ имеем
$I(a)=\int_0^{\pi/2}\ln(a^2)\dx=\pi\ln a$, откуда $\pi\ln(2a)+C=\pi\ln a$ и $C=-\pi\ln2$. Итог:
\[
  I(a)=\pi\ln(a+b)-\pi\ln2=\pi\ln\frac{a+b}{2}.
\]
\end{example}

\begin{example}[применение теоремы Фубини]
$\displaystyle\int_0^1\frac{x^b-x^a}{\ln x}\dx$ ($b>a>0$). Так как
$\dfrac{x^b-x^a}{\ln x}=\int_a^b x^y\,dy$, по теореме Фубини
\[
  \int_0^1\frac{x^b-x^a}{\ln x}\dx=\int_a^b\!\Bigl(\int_0^1 x^y\dx\Bigr)dy
   =\int_a^b\frac{dy}{y+1}=\ln\frac{b+1}{a+1}.
\]
\end{example}

\subsection{Правило Лейбница для переменных пределов}

\begin{theorem}
Если $I(y)=\displaystyle\int_{\alpha(y)}^{\beta(y)}f(x,y)\dx$, причём $f,f'_y$ непрерывны, а
$\alpha,\beta$ дифференцируемы, то
\[
  I'(y)=f\bigl(\beta(y),y\bigr)\beta'(y)-f\bigl(\alpha(y),y\bigr)\alpha'(y)
   +\int_{\alpha(y)}^{\beta(y)}f'_y(x,y)\dx .
\]
\end{theorem}

\begin{example}
$\displaystyle I(y)=\int_{y^2}^{y^3}\frac{e^{xy}}{x}\dx$, тогда
$\displaystyle I'(y)=\frac{e^{y^4}}{y^3}\cdot3y^2-\frac{e^{y^3}}{y^2}\cdot2y
+\int_{y^2}^{y^3}e^{xy}\dx=2\frac{e^{y^4}}{y}-3\frac{e^{y^3}}{y}$.
\end{example}

\section{Интеграл Пуассона}

\begin{theorem}
$\displaystyle\int_0^{+\infty}e^{-x^2}\dx=\frac{\sqrt\pi}{2}$ (и потому
$\int_{-\infty}^{+\infty}e^{-x^2}\dx=\sqrt\pi$).
\end{theorem}

\begin{proof}[Через параметр]
Положим $\displaystyle F(y)=\Bigl(\int_0^y e^{-x^2}\dx\Bigr)^2$. Тогда
\[
  F'(y)=2e^{-y^2}\!\int_0^y e^{-x^2}\dx\overset{x=ty}{=}2\int_0^1 e^{-y^2(t^2+1)}y\,dt .
\]
Введём $G(y)=\int_0^1\dfrac{e^{-y^2(t^2+1)}}{-(t^2+1)}\,dt$; прямой подсчёт даёт $G'(y)=F'(y)$,
значит $F(y)=G(y)+\Const$. При $y=0$: $F(0)=0$, $G(0)=-\int_0^1\frac{dt}{t^2+1}=-\frac\pi4$,
поэтому $\Const=\frac\pi4$. При $y\to+\infty$ имеем $G(y)\to0$, и
\[
  \Bigl(\int_0^{+\infty}e^{-x^2}\dx\Bigr)^2=\lim_{y\to+\infty}F(y)=\frac\pi4
  \ \Rightarrow\ \int_0^{+\infty}e^{-x^2}\dx=\frac{\sqrt\pi}{2}.
\]
\end{proof}

\section{Эйлеровы интегралы}

\subsection{Гамма-функция}

\begin{definition}
При $s>0$ полагают $\displaystyle\Gamma(s)=\int_0^{+\infty}x^{s-1}e^{-x}\dx$.
\end{definition}

Интеграл сходится: на $+\infty$ — за счёт множителя $e^{-x}$, в нуле — поскольку $x^{s-1}$
интегрируема при $s>0$.

\begin{property}\leavevmode
\begin{enumerate}[label=\arabic*),leftmargin=*]
  \item $\Gamma(s+1)=s\,\Gamma(s)$;
  \item $\Gamma(1)=1$, откуда $\Gamma(n+1)=n!$ для $n\in\N_0$;
  \item $\Gamma\!\left(\tfrac12\right)=\sqrt\pi$ (сведением к интегралу Пуассона);
  \item формула дополнения Эйлера: $\Gamma(s)\Gamma(1-s)=\dfrac{\pi}{\sin(\pi s)}$, $0<s<1$.
\end{enumerate}
\end{property}

\begin{proof}[Свойство 1]
Интегрированием по частям ($u=x^s$, $dv=e^{-x}\dx$):
\[
  \Gamma(s+1)=\int_0^{+\infty}x^s e^{-x}\dx
   =\underbrace{-x^s e^{-x}\Big|_0^{+\infty}}_{=0}+s\int_0^{+\infty}x^{s-1}e^{-x}\dx=s\,\Gamma(s).
\]
А $\Gamma(1)=\int_0^{+\infty}e^{-x}\dx=1$, поэтому $\Gamma(n+1)=n\,\Gamma(n)=\dots=n!$.
\end{proof}

\subsection{Бета-функция}

\begin{definition}
При $p,q>0$ полагают $\displaystyle B(p,q)=\int_0^1 x^{p-1}(1-x)^{q-1}\dx$.
\end{definition}

\begin{property}\leavevmode
\begin{enumerate}[label=\arabic*),leftmargin=*]
  \item симметрия: $B(p,q)=B(q,p)$;
  \item тригонометрическая форма (подстановка $x=\sin^2\varphi$):
        $\displaystyle B(p,q)=2\int_0^{\pi/2}\sin^{2p-1}\varphi\,\cos^{2q-1}\varphi\,d\varphi$;
  \item связь с гамма-функцией: $\displaystyle B(p,q)=\frac{\Gamma(p)\,\Gamma(q)}{\Gamma(p+q)}$.
\end{enumerate}
\end{property}

\end{document}
